\documentclass[a4paper, 11pt]{article}

% --- PACCHETTI ---
\usepackage[utf8]{inputenc}
\usepackage[italian]{babel}
\usepackage[margin=2.5cm]{geometry}
\usepackage{graphicx}
\usepackage{tabularx}
\usepackage{hyperref}
\usepackage{fancyhdr}
\usepackage{array}
\usepackage{titlesec}
\usepackage{xcolor}

% --- CONFIGURAZIONE HEADER E FOOTER ---
\pagestyle{fancy}
\fancyhf{}
\fancyhead[L]{Bit Crusher}
\fancyhead[R]{Verbale Interno N. 2} 
\fancyfoot[C]{\thepage}
\renewcommand{\headrulewidth}{0.4pt}

% Configurazione link indice
\hypersetup{
    colorlinks=true,
    linkcolor=black,
    urlcolor=blue
}

\begin{document}

% ==========================================
% FRONTESPIZIO
% ==========================================
\begin{titlepage}
    \centering
    
    % --- LOGHI ---
    % Sostituisci i \framebox con i comandi \includegraphics corretti
    \begin{minipage}{0.45\textwidth}
        \centering
        % Esempio: \includegraphics[width=0.7\textwidth]{logo_unipd.png}
        \includegraphics[width=0.7\textwidth]{unipd_logo.png} 
    \end{minipage}
    \hfill
    \begin{minipage}{0.45\textwidth}
        \centering
        % Esempio: \includegraphics[width=0.7\textwidth]{WhatsApp Image 2026-10-15 at 17.29.12.jpeg}
        \includegraphics[width=0.7\textwidth]{Logo-BitCrusher-Nero.png}
    \end{minipage}

    \vspace{1.5cm}
    {\Large \textbf{UNIVERSITÀ DEGLI STUDI DI PADOVA}} \\
    \vspace{0.3cm}
    {\large Corso di Laurea in Informatica} \\
    \vspace{0.3cm}
    {\large Insegnamento di Ingegneria del Software} \\
    \vspace{0.3cm}
    {\large Anno Accademico 2026/2027} \\

    \vspace{2cm}
    {\Huge \textbf{Verbale di Riunione}} \\
    \vspace{0.5cm}
    {\Large Riunione Interna \#2}

    \vspace{2cm}
    \begin{flushleft}
        \large
        \textbf{Gruppo:} Bit Crusher \\
        \textbf{Email:} bitcrusherteam@gmail.com \\
        \textbf{Versione:} v1.0.0 \\
        \textbf{Data della riunione:} 2026-10-21 \\
        \textbf{Data ultima modifica:} 22 ottobre 2026
    \end{flushleft}
    \vfill
\end{titlepage}

\newpage

% ==========================================
% REGISTRO DELLE MODIFICHE
% ==========================================
\section*{Registro delle Modifiche}
\vspace{0.5cm}
\begin{center}
\renewcommand{\arraystretch}{1.5}
\begin{tabularx}{\textwidth}{|l|l|l|l|X|}
\hline
\textbf{Versione} & \textbf{Data} & \textbf{Autore} & \textbf{Verificatore} & \textbf{Descrizione} \\ \hline
1.0.0 & 2026-10-21 & Antonio Tufano & Riccardo Pozzato & Prima stesura delle sezioni: OdG, Resoconto, Decisioni e To-Do. \\ \hline

\end{tabularx}
\end{center}

\newpage

% ==========================================
% INDICE
% ==========================================
\tableofcontents
\newpage
% ==========================================
% 1. INFORMAZIONI GENERALI
% ==========================================
\section{Informazioni Generali}
\begin{itemize}
    \item \textbf{Luogo/Piattaforma:} Riunione Discord
    \item \textbf{Data:} 2026-10-21
    \item \textbf{Orario di inizio:} 17:30
    \item \textbf{Orario di fine:} 19:45
    \item \textbf{Responsabile:} -
    \item \textbf{Scriba:} Riccardo Pozzato
\end{itemize}

\vspace{0.5cm}
\textbf{Partecipanti presenti:} \\
\vspace{0.2cm}
\renewcommand{\arraystretch}{1.2}
\begin{tabularx}{0.8\textwidth}{|X|c|}
\hline
\textbf{Nome e Cognome} & \textbf{Matricola} \\ \hline
Pietro Leonardo Acampora & 2146762 \\ \hline
Marco Bollettin & 2147955 \\ \hline
Alexandru Mitu & 2101083 \\ \hline
Riccardo Pozzato & 2147988 \\ \hline
Alex Singh & 2148621 \\ \hline
Antonio Tufano & 2137972 \\ \hline
Erik Giuliano & 2138012 \\ \hline
\end{tabularx}

\vspace{0.5cm}
\textbf{Partecipanti assenti:} \\
\vspace{0.2cm}
\renewcommand{\arraystretch}{1.2}
\begin{tabularx}{0.8\textwidth}{|X|c|}
\hline
\textbf{Nome e Cognome} & \textbf{Matricola} \\ \hline
 Nessuno & - \\ \hline
\end{tabularx}

% ==========================================
% 2. ORDINE DEL GIORNO
% ==========================================
\vspace{0.5cm}
\section{Ordine del Giorno}
Gli argomenti previsti per la riunione odierna sono:
\begin{enumerate}
    \item Stesura delle domande per le aziende committenti
    \item Riorganizzazione della repository GitHub e accordo sulle convenzioni per l'utilizzo
    \item Invio mail di decisione nome e logo del gruppo al prof. Vardanega
    \item Caricare nel repository layout LaTeX del verbale e verbale interno n.1
    \item Invio delle mail alle aziende con i relativi quesiti
\end{enumerate}

% ==========================================
% 3. RESOCONTO DELLA RIUNIONE
% ==========================================
\newpage
\section{Resoconto della Riunione}

\subsection{Stesura domande relative ai capitolati} 
Sono state discusse le domande sorte ai membri del gruppo riguardo ogni capitolato proposto e si è valutato quali fosse necessario chiedere alle aziende proponenti.


\subsection{Riorganizzazione repository GitHub}
Il repository GitHub è stato riorganizzato in sezioni, così da avere già un'organizzazione basilare. Inoltre sono state discusse e scelte le convenzioni per l'utilizzo di tale repository.

\subsection{Invio mail al prof. Vardanega}
Inviata al prof. Vardanega la mail per comunicare il nome, il logo e il recapito di posta elettronica scelti dal gruppo, come richiesto.

\subsection{Caricato layout LaTeX e primo verbale nel repository}
Caricati nel repository GitHub i primi documenti interni, il layout LaTeX per i verbali e il file pdf del primo verbale della riunione del 17/10/2026 completo.

\subsection{Invio mail alle aziende}
Sono state mandate le prime mail alle aziende proponenti, dove si chiedono chiarimenti e spiegazioni riguardo a dubbi sorti relativamente ai capitolati. Siamo in attesa di risposta, per poter organizzare le prime riunioni.

% ==========================================
% 4. DECISIONI PRESE
% ==========================================
\section{Decisioni Prese}
\renewcommand{\arraystretch}{1.5}
\begin{tabularx}{\textwidth}{|c|X|}
\hline
\textbf{Codice} & \textbf{Decisione} \\ \hline
V-1.1 & Scelte le domande relative ai capitolati da rivolgere alle aziende \\ \hline
V-1.2 & Scelte le convenzioni per l'utilizzo del repository GitHub\\ \hline


\end{tabularx}

% ==========================================
% 5. AZIONI DA SVOLGERE (TO-DO)
% ==========================================
\vspace{0.5cm}
\section{Azioni da Svolgere (To-Do)}
\begin{tabularx}{\textwidth}{|c|X|l|l|}
\hline
\textbf{Codice Rif.} & \textbf{Descrizione Task} & \textbf{Assegnatario} & \textbf{Scadenza} \\ \hline
VI-1.1 & Continuare il dialogo con le aziende e organizzare le riunioni & Marco Bollettin & 2026-10-28 \\ \hline


\end{tabularx}

% ==========================================
% 6. PROSSIMA RIUNIONE
% ==========================================
\vspace{0.5cm}
\section{Prossima Riunione}
\begin{itemize}
    \item \textbf{Data:}  Da concordare
    \item \textbf{Ora:} Da concordare
    \item \textbf{OdG preliminare:} Discutere sulle riunioni fatte con le aziende, inizio stesura documenti per la candidatura
\end{itemize}

\end{document}